\begin{figure}[!ht]
\centering
\begin{tikzpicture}[x=1cm,y=1cm,>=Latex,font=\small]
 \fill[blue!8] (2.8,.83) rectangle (8.1,1.17);
 \fill[blue!8] (2.8,-1.17) rectangle (8.1,-.83);
 \draw[->] (.6,1)--(9,1);
 \draw[->] (.6,-1)--(9,-1);
 \node[anchor=east] at (.45,1) {physical};
 \node[anchor=east] at (.45,-1) {rescaled};
 \foreach \x in {1.2,2.8,5.3,8.1}{
   \draw (\x,.91)--(\x,1.09);
   \draw (\x,-1.09)--(\x,-.91);
 }
 \node[above] at (1.2,1.1) {$r_0$};
 \node[above] at (2.8,1.1) {$r_{j_\ell}$};
 \node[above] at (5.3,1.1) {$s_\ell$};
 \node[above] at (8.1,1.1) {$q_\ell s_\ell$};
 \node[below] at (1.2,-1.1) {$r_0/s_\ell$};
 \node[below] at (2.8,-1.1) {$a_\ell\asymp q_\ell^{-1}$};
 \node[below] at (5.3,-1.1) {$1$};
 \node[below] at (8.1,-1.1) {$q_\ell$};
 \draw[->,gray] (5.3,.65)--(5.3,-.6)
   node[midway,right,black] {divide radii by $s_\ell$};
\end{tikzpicture}
\caption{The physical and rescaled radii in
Proposition~\ref{scales:good} (schematic, not to scale).
The same observation index supplies both the comparison on the shaded
annulus and the intermediate subsolution.  Its rescaled inner radius
tends to zero while the outer endpoint tends to infinity.}
\label{fig:scales}
\end{figure}
